\documentclass[11pt,a4paper]{article}
\usepackage[margin=1in]{geometry}
\usepackage{amsmath,amssymb,amsthm}
\usepackage{mathtools}
\usepackage{microtype}
\usepackage{enumitem}
\usepackage{xcolor}
\usepackage[colorlinks=true,linkcolor=blue!50!black,urlcolor=blue!50!black]{hyperref}
\usepackage{titlesec}

\titleformat{\section}{\normalfont\large\bfseries}{\thesection.}{0.6em}{}
\newcommand{\ind}[1]{\mathbf{1}\!\left[#1\right]}
\newcommand{\sstar}{s^{\star}}
\newtheorem{proposition}{Proposition}
\theoremstyle{definition}
\newtheorem{remark}{Remark}

\title{\bfseries A scheduled-time feature from projected scores\\[2pt]
\large occupancy fractions, the weighted survival function, and where the kernel goes}
\author{}
\date{20 August 2026}

\begin{document}
\maketitle
\vspace{-2.5em}

\section*{The question, as posed}
\noindent Use the mining / stripping capacity to define the fraction of a time period that mining
each block would take, for the sequence of blocks passed to the transformer. Given those
fractions, build a feature $t$ that respects precedence in the projected score, following the rule
\[
  t_j \;=\; \tau_j + t_i \qquad\text{if}\qquad \sstar(x_i) > \sstar(x_j),
\]
with $i,j$ block indices and $\sstar$ the projected score. Consider whether the kernel
interpolation technique can build $t$ from the projected scores.

\medskip
\noindent\textit{Grounding note.} \texttt{soft\_rank.py:soft\_npv} already computes
$W_r[i]=\sum_j w_{r,j}\,\sigma\!\big((s_j-s_i)/\tau\big)$, the resource consumed before block $i$.
That is the object below, up to a division by capacity. What is new is exposing it as a
\emph{feature} and building it in $O(n\log n)$ rather than $O(n^2)$.

\section{The occupancy fraction}

Let block $j$ have consumption $w_{r,j}$ on resource $r$ and let $C_r$ be the per-period capacity
of that resource. In \texttt{capacity\_cuts.make\_resource} the axes are already
\[
  \text{mining} = \text{tonnage},\qquad
  \text{processing} = \text{tonnage}\cdot\ind{\text{ore}},\qquad
  \text{stripping} = \text{tonnage}\cdot\ind{\text{waste}} .
\]
The fraction of a period that block $j$ occupies is
\[
  \boxed{\;\tau_j \;=\; \max_r \frac{w_{r,j}}{C_r}\;}
\]
The maximum, not the mining term alone: a period ends when the \emph{first} resource binds, so
charging only mining tonnage silently assumes the mill never binds.

Note that $\tau$ is not the stripping ratio. The ratio $w_{\mathrm{waste},j}/w_{\mathrm{ore},j}$
enters only through \emph{which} axis attains the maximum. A separate stripping cap is redundant
unless waste handling is independently constrained, since
$\text{mining}=\text{processing}+\text{stripping}$ leaves only two of the three axes free.

With time measured in periods, $\sum_j \tau_j$ is the horizon needed to mine everything, and a
$T$-period schedule of constant capacity is feasible only if $\sum_j \tau_j \le T$.

\section{The recursion telescopes into a weighted survival function}

The rule $t_j = \tau_j + t_i$ is correct only when $i$ is the \emph{immediate} predecessor of $j$
in score order. Applied to every higher-scoring $i$ it double-counts. Write $\pi(j)$ for the block
whose score is the smallest one exceeding $\sstar_j$; then $t_j = \tau_j + t_{\pi(j)}$ telescopes
along the score order to the closed form
\[
  \boxed{\;
  t_j \;=\; \tau_j \;+\; G\!\left(\sstar_j\right),
  \qquad
  G(s) \;\coloneqq\; \sum_{i} \tau_i \,\ind{\sstar_i > s}
  \;}
\]
So $t_j$ is not a pairwise object at all. It is the block's own occupancy plus the
\textbf{$\tau$-weighted survival function of the score distribution, evaluated at the block's own
score}. Equivalently, with $\mu = \sum_i \tau_i \delta_{\sstar_i}$ the score-mass measure,
$G(s) = \mu\big((s,\infty)\big)$.

$G$ is scalar-in, scalar-out, non-increasing, and piecewise constant with $n$ jumps of height
$\tau_i$. Everything cheap below follows from that one fact. Here $G(\sstar_j)$ is the
\emph{start} time of block $j$ and $t_j$ its \emph{completion} time.

\begin{proposition}[precedence is inherited, at no cost]
If $(i \to j) \in E$ and the projection has delivered $\sstar_i \ge \sstar_j$, then
$G(\sstar_i) \le G(\sstar_j)$, i.e.\ block $j$ never \emph{starts} before block $i$.
\end{proposition}
\begin{proof}
$G$ is non-increasing and $\sstar_i \ge \sstar_j$.
\end{proof}

\noindent No extra constraint is required: precedence in score space transfers to the time feature
for free. But the inheritance is \emph{weak}. Exact ties give equal start times, which the
resource constraint cannot honour. Ties therefore need breaking, either lexicographically
($+\varepsilon\cdot\mathrm{rank}$) or by splitting the tied mass,
\[
  t_j \;=\; \tau_j + G(\sstar_j) + \tfrac12 \sum_i \tau_i\, \ind{\sstar_i = \sstar_j}.
\]

\section{Kernel interpolation --- yes, but in score space}

Smooth the step function with a kernel CDF $\Phi$ of bandwidth $h$:
\[
  G_h(s) \;=\; \sum_i \tau_i \,\Phi\!\left(\frac{\sstar_i - s}{h}\right).
\]
With $\Phi$ logistic this is \emph{identically} the pairwise sigmoid already inside
\texttt{soft\_npv}. The soft-rank you have and the kernel-smoothed survival function are the same
operator seen from two sides. The existing code is $O(n^2)$ only because it evaluates $G_h$
separately at each $\sstar_j$, when $G_h$ is a single univariate curve.

The interpolation version:
\begin{enumerate}[leftmargin=1.4em,itemsep=1pt]
  \item choose $m \ll n$ anchor knots $\sigma_1 > \sigma_2 > \dots > \sigma_m$ at
        $\tau$-weighted score quantiles, so knots are dense where score mass is dense;
  \item evaluate $G_h$ exactly at the knots --- $O(nm)$, or $O(n\log n)$ for the hard $G$ by
        sorting the scores and taking a cumulative sum of $\tau$;
  \item recover every $t_j$ by interpolating $G$ at $\sstar_j$.
\end{enumerate}

\begin{remark}[the interpolant must be monotone]
Use PCHIP, or linear interpolation on the monotone knot grid. A plain cubic spline and a GP
posterior mean can both overshoot, and an overshoot destroys the implication in Proposition~1 ---
the whole point of the construction.
\end{remark}

Cost falls from $O(n^2)$ to $O(n\log n + nm)$, which is the difference between a feature that can
be recomputed on every forward pass and one that cannot.

\begin{remark}[this is a different kernel from \texttt{anchor\_interpolation.py}]
There, the kernel lives in augmented block-geometry space $[z_i;\lambda c_i]$ and its job is to
spread a constraint correction laterally across the deposit. Here it lives in one-dimensional
score space and its job is to smooth a monotone reduction. Same word, unrelated operators. Do not
reuse \texttt{build\_features} or the Wendland radius for this.
\end{remark}

\section{What the feature buys, and the gradient it carries}

Once $t_j$ exists, everything downstream is differentiable in closed form. The period index
$p_j = \lceil t_j \rceil$ softens to
\[
  p_j = \sum_k \sigma\!\left(\frac{t_j - k}{\varepsilon}\right),
  \qquad
  \pi_{jk} = \sigma\!\left(\frac{t_j - k + 1}{\varepsilon}\right) - \sigma\!\left(\frac{t_j-k}{\varepsilon}\right),
\]
the latter a soft one-hot over periods; and the discount is simply $\gamma^{\,t_j}$, with no
capacity-crossing sigmoid stack, because $t$ is \emph{already} in period units.

The gradient is the informative part. With $\phi = \Phi'$,
\[
  \frac{\partial t_j}{\partial \sstar_i}
  \;=\; -\,\frac{\tau_i}{h}\,\phi\!\left(\frac{\sstar_i - \sstar_j}{h}\right)
  \quad (i \ne j),
  \qquad
  \frac{\partial t_j}{\partial \sstar_j}
  \;=\; \frac{1}{h}\sum_i \tau_i\,\phi\!\left(\frac{\sstar_i - \sstar_j}{h}\right)
  \;=\; \hat f_\tau\!\left(\sstar_j\right),
\]
the $\tau$-weighted kernel density of scores at $\sstar_j$. Two consequences.

\medskip
\noindent\textbf{(i) $h$ is a receptive field in score space}, not a generic temperature. Only
blocks within $O(h)$ of $j$ \emph{in score} exchange gradient. As $h \to 0$ the schedule becomes
exact but the gradient collapses onto near-ties and training stalls; large $h$ gives a dense but
badly biased $t$. Anneal $h$ as a quantile spacing --- e.g.\ the score gap at the $q$-th
percentile --- so it stays scale-free as the score distribution drifts during training.

\medskip
\noindent\textbf{(ii) Gradient magnitude tracks local score density.} A block in a crowded score
region moves its own $t$ hard; a score outlier is nearly frozen. If that shows up as pathological,
normalise by $\hat f_\tau$ and accept the resulting bias.

\section{The seam to settle before building}

$t$ is a function of $\sstar$, which is the network's output. Feeding $t$ back in as an input
feature is therefore a fixed-point equation $\sstar = F\big(x,\, t(\sstar)\big)$, not a feature
map. Three coherent readings, and they are materially different builds:

\begin{description}[leftmargin=2.2em,style=nextline,itemsep=3pt]
  \item[(a) loss-side only.] $t$ never enters the encoder; replace the sigmoid stack in
    \texttt{soft\_npv} with $\gamma^{\,t_j}$ and stop there.
  \item[(b) decoder feature, one unrolled pass.]
    $s^{(1)} = F(x)$, then $t = \tau + G_h\!\big(s^{(1)}\big)$, then $s^{(2)} = F(x, t)$,
    backpropagating through both passes. This is what ``a feature based on the projected score''
    most naturally means.
  \item[(c) target-BO seam.] $t$ enters \texttt{target\_bo}'s \texttt{soft\_profile} as a
    scheduled-time target rather than a score target.
\end{description}

\medskip
\noindent My reading of the request is \textbf{(b)}. Three things fix the operator completely, and
I want them agreed in writing before code exists:

\begin{enumerate}[leftmargin=1.6em,itemsep=2pt]
  \item which of (a)/(b)/(c);
  \item the tie-break rule --- lexicographic or split-mass;
  \item whether $\tau_j = \max_r w_{r,j}/C_r$ over all resources, or mining tonnage alone.
\end{enumerate}

\noindent The last two times this operator was built from an unconfirmed reading, it was the wrong
operator.

\end{document}
